% Бүлэг 1

\chapter{Удиртгал}
\label{Chapter1}
\pagecolor{white}

\newcommand{\keyword}[1]{\textbf{#1}}
\newcommand{\tabhead}[1]{\textbf{#1}}
\newcommand{\code}[1]{\texttt{#1}}
\newcommand{\file}[1]{\texttt{\bfseries#1}}
\newcommand{\option}[1]{\texttt{\itshape#1}}

\section{Сэдэв сонгосон үндэслэл}

Орчин үеийн дижитал орчинд богино хэмжээний видео контент нь нийгмийн сүлжээнд давамгайлж буй гол хэлбэр болж хувирчээ. TikTok, YouTube Shorts, Instagram Reels зэрэг платформууд нь сар бүр тэрбум гаруй хэрэглэгчид хүрч байгаа бөгөөд уг платформуудад агуулга бүтээх эрэлт нь маш өндөр байна.

Чанартай видео контент бүтээх нь цаг хугацаа, мэдлэг, хэрэгсэл зэрэг олон нөөц шаарддаг. Гараар видео хийхэд:

\begin{itemize}
\item Скрипт бичих --- 1-3 цаг
\item Зураг/дизайн хийх --- 2-5 цаг
\item Дуу бичлэг хийх --- 1-2 цаг
\item Видео монтаж хийх --- 3-8 цаг
\end{itemize}

Нийт нэг видеод 1-3 өдөр зарцуулагддаг. AI технологи ашиглан эдгээр үйл ажиллагааг автоматжуулснаар нэг видеог 5-15 минутад бэлдэх боломж бий болно.

Монгол хэлний видео контент интернэтэд харьцангуй бага байдаг бөгөөд Монгол компани, брэнд, блогчдод зориулсан хялбар, автомат видео контент үүсгэгч байхгүй байна. Олон улсын зах зээлд InVideo AI, Pictory, Synthesia, Runway Gen-3 зэрэг системүүд ажиллаж байгаа боловч аль ч систем Монгол хэлийг дэмждэггүй байна. Энэ нь манай системийн зах зээлийн том цоорхойг харуулж байна.

Иймд зөвхөн контент автоматаар үүсгэхэд зогсохгүй --- үүсгэсэн видеогоо AI Video Studio засварлагчаар scene бүрийн narration, зураг, дуу хоолой, background music, subtitle засварлаж боловсруулах боломж бүхий цогц систем хөгжүүлэх шаардлага тодорхой болов.

\section{Зорилго}

Энэхүү дипломын ажлын гол зорилго нь хэрэглэгч зөвхөн сэдэв, жанр, урт зэрэг үндсэн тохиргоог оруулахад дараах бүх үйл явцыг автоматаар гүйцэтгэж, мөн үүсгэсэн контентоо Studio засварлагчаар боловсруулах боломж бүхий хиймэл оюунд суурилсан цогц видео контент үүсгэлтийн веб систем хөгжүүлэх юм:

\begin{enumerate}
\item Монгол болон Англи хэлний скрипт (үг, дүр зураглал, дуу хоолойн дэлгэрэнгүй) автоматаар бүтээх
\item Scene бүрд тохирсон зураг (AI) үүсгэж, видеоны үндэс болгох
\item Скриптийн хиймэл дуу хоолой синтез хийх (TTS)
\item Зураг, дуу, текстийг нэгтгэн MP4 видео файл гаргах
\item AI Video Studio засварлагчаар scene бүрийн narration, зураг, дуу хоолой, background music, subtitle засварлах
\item Зөвхөн өөрчлөгдсөн scene-үүдийг хэсэгчилсэн re-render хийх
\item Үүлэн хадгалалтад байршуулж, YouTube/TikTok/Instagram хэлбэрт тааруулан хэрэглэгчид хүргэх
\end{enumerate}

\section{Зорилт}

Дипломын ажлын хүрээнд дараах зорилтуудыг дэвшүүлж байна:

\begin{enumerate}
\item Монгол болон олон улсын ижил төстэй системүүдийг судалж, харьцуулсан шинжилгээ хийх;
\item Ашиглах AI загварууд (Groq, Google Gemini, ElevenLabs, Magic Hour) болон технологийн судалгаа хийх;
\item Системийн хэрэглэгчийн болон функциональ шаардлагуудыг тодорхойлж, Use Case диаграм гаргах;
\item AI Video Studio засварлагчийн Use Case, шаардлага, урсгалыг тодорхойлох;
\item PostgreSQL өгөгдлийн сангийн нормчилсон схем, хүснэгтийн бүтцийг зохиомжлох;
\item Role-Based Access Control (RBAC) бүхий JWT баталгаажуулалтын тогтолцоог хэрэгжүүлэх;
\item Node.js + Express + TypeScript ашиглан RESTful API backend хөгжүүлэх;
\item FFmpeg ашиглан зураг, дуу, текстийг нэгтгэн видео гаргах болон хэсэгчилсэн re-render модуль хөгжүүлэх;
\item Cloudinary болон AWS S3-д видеог байршуулах, хэрэглэгчид хүргэх модуль хэрэгжүүлэх.
\end{enumerate}

\section{Бүлгийн дүгнэлт}

Энэ бүлэгт Монгол хэлний видео контентын хомсдол болон одоо байгаа системүүдийн дүн шинжилгээнд үндэслэн сэдвийн шаардлага болон зорилтуудыг тодорхойлсон. Харьцуулсан шинжилгээгээр тус дипломын ажил нь Монгол хэлний дэмжлэг, AI Video Studio засварлагч, хэсэгчилсэн re-render, олон платформ export зэрэг одоо байгаа системүүдэд байхгүй шинэлэг боломжуудыг нэгтгэсэн цогц шийдэл болохыг тодорхойлов.
